(1) 左侧波形 \(y_L=e^{ikx-i\omega t}+re^{-ikx-i\omega t}\)，右侧为 \(y_R=te^{ikx-i\omega t}\)。弦在质量处不断裂，故位移连续给 \(1+r=t\)。小珠受左右张力的净横向力为 \(T[y_x(0^+)-y_x(0^-)]\)，其牛顿方程是 \(T[y_x(0^+)-y_x(0^-)]=-M\omega^2y(0)\)。跳跃方向由两端拉力指向反向决定，不应把点质量当成固定端。

(2) 导数跳跃为 \(ikT[t-(1-r)]=2ikTr\)，所以 \(2ikTr=-M\omega^2t=-M\omega^2(1+r)\)。令 \(\zeta=M\omega/(2T k/\omega)=M\omega/(2\sqrt{T\mu})\)；更直接由 \(2Tk=M\omega^2/\zeta\)，故 \(r=-i\zeta/(1+i\zeta)\)，\(t=1/(1+i\zeta)\) 在选取 \(e^{-i\omega t}\) 时可通过原式代入核验：\(2i r=-2\zeta t\) 推出 \(r=i\zeta/(1-i\zeta)\)，\(t=1/(1-i\zeta)\)。后者是与所给跳跃符号严格一致的表达，取定上述跳跃符号后，不能再随意改变散射相位的正负号。

(3) 采用由跳跃方程确定的 \(r=i\zeta/(1-i\zeta)\)、\(t=1/(1-i\zeta)\)，有 \(|r|^2=\zeta^2/(1+\zeta^2)\)、\(|t|^2=1/(1+\zeta^2)\)，二者和为一。因左右弦张力及线密度相同，能流权重相同，这就是整个无耗散点质量散射的功率守恒。若 \(\omega\to0\)，\(\zeta\to0\) 则反射趋零；\(\omega\to\infty\) 时 \(|r|\to1\) 且透射趋零。

(4) 由于 \(\zeta=M\omega/(2\sqrt{T\mu})\)，质量越大或频率越高，对弦的惯性障碍越强，点处透射振幅大小 \(1/\sqrt{1+\zeta^2}\) 越低，反射相位也由复杂阻抗决定。若 \(M\to0\)，\(r\to0,t\to1\)，弦退化为完全均匀无缺陷传播。该点质量模型没有自身弹簧，因此不具有限频孤立局部共振；若观察到窄线共振必须增加局部回复力或其他物理结构。
